\chapter*{Preface}
\addcontentsline{toc}{chapter}{Preface}

The large language model revolution has moved faster than any previous technology cycle in machine learning. What required a year-long research effort at a top lab in 2022 now fits in a weekend hackathon; models that once cost \$10 million to train can be replicated for tens of thousands of dollars.

This guide is written for engineers and researchers who need to \emph{do the work}---not survey it from a distance. It covers the full training stack: data pipelines, pre-training at scale, the post-training alignment stack, evaluation, safety, and production deployment. Each chapter synthesizes the most important findings from 2024--2026 alongside the enduring principles from earlier years.

\textbf{Who this book is for:} ML engineers building fine-tuning pipelines; researchers designing new alignment methods; infrastructure engineers provisioning GPU clusters; technical leads making architectural decisions for LLM-powered products.

\textbf{How to read this book:} This guide is designed for both linear reading and targeted reference.
\begin{itemize}
  \item \textbf{Parts I--III (Foundations \& Pre-training):} Essential for teams building or modifying base models. Covers data, architecture, and scaling.
  \item \textbf{Parts IV--V (Post-training \& Alignment):} The high-value stack for 2026. Relevant to everyone from hobbyists to frontier researchers.
  \item \textbf{Part VI (Advanced Topics):} Specialized research on reasoning, multi-modality, and the 2027 roadmap.
\end{itemize}

Each chapter begins with a \textbf{What You Will Learn} overview and a \textbf{Mini Table of Contents} for quick navigation. Mathematical formulations are concentrated in Appendix~\ref{app:pipeline}--\ref{app:continual}, with direct links provided throughout the text. Ready-to-use PyTorch code exists in Appendix~H, serving as the executable companion to the theoretical chapters.
